% !TEX root = ../math.tex
\section{Spines and Pair Canonicality}
\label{sec:bpe-spines}

For implementation purposes it is useful to represent a BPE merge rule as a tuple
\[
(\iota_1,\iota_2) \mapsto (\iota_3, r),
\]
where $\iota_1$ and $\iota_2$ are the IDs of the two child tokens, $\iota_3$ is the ID of
their merged token, and $r$ is the merge rank. Lower rank means higher priority.

This gives three natural lookup views of the same merge:
\[
(\iota_1,\iota_2) \mapsto (\iota_3,r),
\qquad
(\iota_1,\iota_3) \mapsto (\iota_2,r),
\qquad
(\iota_2,\iota_3) \mapsto (\iota_1,r).
\]
The first view answers the usual forward question, while the latter two answer
``one child plus the parent'' queries that arise when reconstructing derivations.

\begin{definition}[Left and right spines]
    Let $\boldsymbol{\tau}$ be a BPE token, and let its binary merge tree be the unique tree whose
    leaves are characters and whose internal nodes are merge results, with root $\boldsymbol{\tau}$.
    The \emph{left spine} of $\boldsymbol{\tau}$ is the leftmost root-to-leaf path, reported from leaf to root.
    The \emph{right spine} is defined symmetrically.
\end{definition}

In the implementation we store either spine as a list of pairs
\[
[(\iota_0, r_0), (\iota_1, r_1), \ldots, (\iota_m, r_m)],
\]
where:
\begin{enumerate}
    \item $\iota_0$ is the ID of the boundary-adjacent character on that side;
    \item $\iota_m$ is the ID of the whole token;
    \item for $j < m$, the rank $r_j$ is the merge rank of the operation that grows $\iota_j$
    into $\iota_{j+1}$ along that spine;
    \item $r_m = \texttt{None}$.
\end{enumerate}

Thus, for a token whose derivation is
\[
h+e \to he,
\qquad
l+l \to ll,
\qquad
he+ll \to hell,
\]
the left spine is
\[
[(h,\operatorname{rank}(h,e)), (he,\operatorname{rank}(he,ll)), (hell,\texttt{None})],
\]
while the right spine is
\[
[(l,\operatorname{rank}(l,l)), (ll,\operatorname{rank}(he,ll)), (hell,\texttt{None})].
\]

\begin{definition}[Spine frontier]
    Let $\boldsymbol{\tau_1}$ and $\boldsymbol{\tau_2}$ be tokens.
    Suppose we are given the right spine of $\boldsymbol{\tau_1}$ and the left spine of
    $\boldsymbol{\tau_2}$, both stored leaf-to-root as above.
    A \emph{spine frontier state} is a pair of indices $(i,j)$ pointing into these two lists.
    The currently exposed boundary tokens are the IDs carried by those two entries.
\end{definition}

\begin{proposition}[Spine-based canonicality algorithm]
    \label{prop:spine-based-canonicality}
    Let $\boldsymbol{\tau_1}$ and $\boldsymbol{\tau_2}$ be BPE tokens.
    Assume their right and left spines are known.
    Consider the following procedure.

    Initialize two pointers at the leaf ends of the spines.
    At each step examine three candidate ranks:
    \begin{enumerate}
        \item the rank of the next merge on the right spine of $\boldsymbol{\tau_1}$;
        \item the rank of the next merge on the left spine of $\boldsymbol{\tau_2}$;
        \item the rank of the merge of the two currently exposed boundary tokens,
        if such a merge exists in the BPE merge table.
    \end{enumerate}
    If the third rank is the unique highest-priority rank, the boundary is crossed and the pair is not canonical.
    If the first or second rank has highest priority, advance the corresponding spine pointer to the next larger
    spine token. If all three ranks are absent, terminate and return that the pair is canonical.

    This procedure returns \texttt{true} if and only if
    \[
    \psi(\boldsymbol{\tau_1},\boldsymbol{\tau_2}).
    \]
\end{proposition}
\begin{proof}
    The key invariant is that after all merges of rank smaller than the current frontier have been applied,
    the exposed token on the left side of the boundary is exactly the current node of the right spine of
    $\boldsymbol{\tau_1}$, and the exposed token on the right side of the boundary is exactly the current node
    of the left spine of $\boldsymbol{\tau_2}$.

    Any merge that is strictly internal to the left token and affects the boundary must enlarge the
    boundary-adjacent suffix of $\boldsymbol{\tau_1}$, and therefore must be the next merge on its right spine.
    Likewise, any boundary-relevant internal merge on the right token must be the next merge on its left spine.
    The only other possible boundary-relevant event is a direct merge of the two currently exposed boundary tokens.
    Therefore every event that can affect whether the boundary survives is represented by exactly one of the
    three candidate ranks listed above.

    If the cross-boundary rank is the highest-priority one, then the next relevant event merges the two boundary
    tokens together, so the boundary does not survive and the pair is not canonical.
    If instead one of the within-token spine ranks is the highest-priority one, then the corresponding token grows
    inward-to-outward along its spine while the boundary is preserved, so advancing that pointer is correct.
    Finally, if all three ranks are absent, no further merge can affect the boundary, so the boundary survives to
    the end of tokenization and the pair is canonical.
\end{proof}

The resulting decision procedure is linear in the two spine lengths:
\[
O(|\operatorname{rspine}(\boldsymbol{\tau_1})| + |\operatorname{lspine}(\boldsymbol{\tau_2})|),
\]
with one merge-table lookup per iteration.
In practice this is substantially cheaper than retokenizing the concatenated string
$\rho(\boldsymbol{\tau_1}) + \rho(\boldsymbol{\tau_2})$ from scratch.
